\subsection{Critérios de Desempenho}
\label{sec:criterios_desempenho}

A avaliação do controle adotado nesta dissertação é em função das trajetórias relativas, à atitude e ao MR.
Os critérios considerados são:
(i) o erro de posição relativa nos eixos $(x,y,z)$ do referencial LVLH;  
(ii) o erro de velocidade relativa;  
(iii) o tempo de acomodação;  
(iv) sobressinal máximo;  
(v) erro de atitude;  
e (vi) a suavidade do esforço de controle.
A escolha do controlador PD é baseado na simplicidade de implementação, baixa exigência computacional e robustez em condições operacionais.
Por outro lado, métodos alternativos, como o controle ótimo linear (LQR) e o controle preditivo baseado em modelo (MPC), podem reduzir os custos em termos de consumo de energia e da otimização de trajetória, mas possuem maior complexidade de projeto e maior custo computacional.
Sendo assim, para os objetivos desta pesquisa, o controlador PD atende os requisitos básicos de rastreamento e estabilização.

As condições de sucesso da missão foram definidas como:  
(i) a convergência da atitude relativa para zero;  
(ii) velocidade relativa inferior ao limite operacional de captura;  
(iii) ausência de violações do corredor de aproximação (\emph{keep-out zone}) e 
(iv) ausência de violações do espaço de trabalho do manipulador robótico.  
